% =============================================================
%  CHAPTER 3 - RELEASE 1 : AUTHENTICATION & USER MANAGEMENT
% =============================================================

\chapteroverview{Release 1 : Authentication \& User Management}{
  \item Introduction
  \item Sprint Backlog
  \item Refined Use Cases
  \item Sequence Diagrams
  \item Class Diagram of the Module
  \item Achieved Increment
  \item Conclusion
}

\section{Introduction}

The first release lays the security foundations of the platform. It delivers everything that revolves around \textit{identity} : registration, secure authentication, profile management and the personal user dashboard. Without this layer, no other module can be safely built on top.

% ---------- 3.2 Sprint backlog ----------
\section{Sprint Backlog}

\begin{table}[H]
\centering
\renewcommand{\arraystretch}{1.35}
\begin{tabularx}{\textwidth}{|>{\columncolor{softblue}}c|X|c|c|}
\hline
\rowcolor{primary}
\textbf{\color{white} ID} & \textbf{\color{white} Task} & \textbf{\color{white} Sprint} & \textbf{\color{white} Status}\\
\hline
T1 & Database design (users, roles, sessions) & S1 & \checkmark \\\hline
T2 & Sign-up form with field validation & S1 & \checkmark \\\hline
T3 & Login with JWT \& password hashing & S1 & \checkmark \\\hline
T4 & Forgot / reset password & S1 & \checkmark \\\hline
T5 & User profile : photo \& avatar upload & S2 & \checkmark \\\hline
T6 & Change password from profile & S2 & \checkmark \\\hline
T7 & Personal dashboard widgets & S2 & \checkmark \\\hline
T8 & Role-based access middleware & S2 & \checkmark \\\hline
\end{tabularx}
\caption{Sprint backlog -- Release 1}
\label{tab:sb1}
\end{table}

% ---------- 3.3 Use cases ----------
\section{Refined Use Cases}

\subsection{Detailed use case : \textit{Sign-up}}

\begin{table}[H]
\centering
\renewcommand{\arraystretch}{1.3}
\begin{tabularx}{\textwidth}{|>{\columncolor{softblue}}p{3.5cm}|X|}
\hline
\textbf{Title}        & Sign-up \\\hline
\textbf{Actor}        & User \\\hline
\textbf{Goal}         & Create a personal account on the platform. \\\hline
\textbf{Pre-condition}& The user is not yet registered. \\\hline
\textbf{Main scenario}& 1. The user opens the sign-up page. \newline
                        2. The system displays the form. \newline
                        3. The user fills in the required fields. \newline
                        4. The system validates the data. \newline
                        5. The system creates the account and sends a confirmation. \\\hline
\textbf{Alternative}  & 4.a If the e-mail already exists, the system displays an error. \\\hline
\textbf{Post-condition}& The account is created and ready to be activated. \\\hline
\end{tabularx}
\caption{Detailed use case : Sign-up}
\label{tab:uc-signup}
\end{table}

\subsection{Use case diagram of Release 1}

\begin{figure}[H]
\centering
\begin{tikzpicture}[
  every node/.style={font=\small},
  uc/.style={ellipse, draw=primary, fill=softblue, minimum width=2.6cm, align=center, inner sep=2pt}
]
  \draw[draw=primary, thick, rounded corners=8pt] (0,-3.5) rectangle (8,3.5);
  \node[font=\bfseries\color{primary}] at (4,3.1) {Authentication \& User Module};

  \node[uc] (su) at (4,2.3) {Sign-up};
  \node[uc] (lo) at (4,1.3) {Login};
  \node[uc] (fp) at (4,0.3) {Forgot password};
  \node[uc] (mp) at (4,-0.7) {Manage profile};
  \node[uc] (cp) at (4,-1.7) {Change password};
  \node[uc] (db) at (4,-2.7) {View dashboard};

  \node[align=center] (user) at (-2,0)
       {\tikz\draw[primary, thick] (0,0) circle (0.18) (0,-0.18) -- (0,-0.7)
         (-0.4,-0.4) -- (0.4,-0.4) (-0.25,-1.05) -- (0,-0.7) -- (0.25,-1.05);\\\textbf{User}};

  \draw (user) -- (su);
  \draw (user) -- (lo);
  \draw (user) -- (fp);
  \draw (user) -- (mp);
  \draw (user) -- (cp);
  \draw (user) -- (db);
\end{tikzpicture}
\caption{Use case diagram -- Release 1}
\label{fig:uc-r1}
\end{figure}

% ---------- 3.4 Sequence diagrams ----------
\section{Sequence Diagrams}

\subsection{Sequence : Login}

\begin{figure}[H]
\centering
\begin{tikzpicture}[
  every node/.style={font=\small},
  lifeline/.style={dashed, color=primary},
  msg/.style={->, thick, color=primary},
  ret/.style={->, thick, dashed, color=secondary},
  actor/.style={rectangle, draw=primary, fill=softblue, minimum width=1.8cm, minimum height=0.7cm}
]
  \node[actor] (u)  at (0,0)  {User};
  \node[actor] (ui) at (3,0)  {UI};
  \node[actor] (api)at (6,0)  {API};
  \node[actor] (db) at (9,0)  {DB};

  \foreach \n in {u,ui,api,db} \draw[lifeline] (\n.south) -- ++(0,-7);

  \draw[msg] (0,-1)  -- node[above,font=\scriptsize]{enter credentials} (3,-1);
  \draw[msg] (3,-2)  -- node[above,font=\scriptsize]{POST /login} (6,-2);
  \draw[msg] (6,-3)  -- node[above,font=\scriptsize]{SELECT user} (9,-3);
  \draw[ret] (9,-4)  -- node[above,font=\scriptsize]{user data} (6,-4);
  \draw[ret] (6,-5)  -- node[above,font=\scriptsize]{200 + JWT token} (3,-5);
  \draw[ret] (3,-6)  -- node[above,font=\scriptsize]{redirect dashboard} (0,-6);
\end{tikzpicture}
\caption{Sequence diagram of the login process}
\label{fig:sd-login}
\end{figure}

% ---------- 3.5 Class diagram ----------
\section{Class Diagram of the Module}

\begin{figure}[H]
\centering
\begin{tikzpicture}[
  every node/.style={font=\small},
  class/.style={rectangle split, rectangle split parts=3, rectangle split part fill={primary, white, white},
                draw=primary, text=white, rectangle split draw splits=true, align=left, inner sep=4pt}
]
  \node[class] (user) at (0,0) {%
    \nodepart{one}\bfseries User
    \nodepart[text=black]{two}
      - id : int\\
      - email : string\\
      - password : string\\
      - firstName : string\\
      - lastName : string\\
      - photo : string\\
      - avatar : string
    \nodepart[text=black]{three}
      + register()\\
      + login()\\
      + updateProfile()\\
      + changePassword()
  };

  \node[class, right=2.5cm of user] (role) {%
    \nodepart{one}\bfseries Role
    \nodepart[text=black]{two}
      - id : int\\
      - name : string\\
      - permissions : json
    \nodepart[text=black]{three}
      + has(perm)
  };

  \node[class, below=1.5cm of user] (session) {%
    \nodepart{one}\bfseries Session
    \nodepart[text=black]{two}
      - id : int\\
      - token : string\\
      - createdAt : date\\
      - expiresAt : date
    \nodepart[text=black]{three}
      + isValid()
  };

  \draw[->, thick, color=primary] (user) -- node[above,font=\scriptsize]{1..*} (role);
  \draw[->, thick, color=primary] (user) -- node[right,font=\scriptsize]{1..*} (session);
\end{tikzpicture}
\caption{Class diagram -- Authentication module}
\label{fig:cd-r1}
\end{figure}

% ---------- 3.6 Achieved increment ----------
\section{Achieved Increment}

At the end of Release 1, the platform delivers :
\begin{itemize}[leftmargin=2em,itemsep=0.2em]
  \item A complete sign-up / login system based on JWT.
  \item Profile management with photo and avatar upload.
  \item Password reset and change.
  \item A first version of the personal dashboard for the user.
  \item A role-based middleware ready to be reused by the next releases.
\end{itemize}

\paragraph{Indicative interfaces.}
The figure below positions, schematically, the screens delivered at the end of this release.

\begin{figure}[H]
\centering
\begin{tikzpicture}[
  scr/.style={rectangle, draw=primary, fill=white, rounded corners=4pt,
              minimum width=3.2cm, minimum height=2.0cm, align=center, font=\small\bfseries\color{primary}}
]
  \node[scr] (a) at (0,0) {Sign-up\\screen};
  \node[scr] (b) at (3.6,0) {Login\\screen};
  \node[scr] (c) at (7.2,0) {Profile\\screen};
  \node[scr] (d) at (10.8,0) {User\\dashboard};
\end{tikzpicture}
\caption{Screens delivered at the end of Release 1}
\label{fig:r1-screens}
\end{figure}

\section{Conclusion}

This first release has stabilized a robust security and identity layer. The next chapter builds on top of it to offer the functional heart of the platform : the training catalog and the course-management module.
